% =====================================================================
% Group C -- Chapter 8: orbit perturbations (one question per exam)
% Source: latex_notes/SFM_08_Orbit Perturbations/cap08.tex
% =====================================================================
\qgroup{C}{Chapter 8 -- Orbit perturbations}

One of the three orbital-mechanics questions is always on Chapter 8. Per the course program, for SRP and drag only the
definition, the expression and the overall effect may be asked (not the derivation of the exact components), and for the
third body the positions of Sun and Moon are not requested.

Perturbed motion: $\ddot{\underline r}=-\frac{\mu}{r^3}\underline r+\underline f$; $\underline h$ and $\underline e$ are no longer constant:
$\dot{\underline h}=\underline r\times\underline f$. Order of magnitude at 1000\,km of altitude ($a_0=\mu/r^2$):
$a_{J2}\simeq10^{-2}a_0$, $a_{3B}\simeq10^{-7}a_0$, $a_{RP}\simeq10^{-9}a_0$, $a_D\simeq10^{-10}a_0$.
Asphericity and third body are conservative; drag and SRP are \emph{surface} perturbations (depend on the spacecraft area);
all of them except SRP depend on the altitude.

\begin{center}
\small
\begin{tabular}{@{}lllr@{}}
\toprule
Q & Topic & Asked & Page\\
\midrule
\ref{q:C1} & Solar radiation pressure: nature and formula & $\times$3, Jan 2023 & \pageref{q:C1}\\
\ref{q:C2} & SRP + geometry of eclipses & Jun 2025 & \pageref{q:C2}\\
\ref{q:C3} & $J_2$: nature, two averaged effects, lat/long dependence & list & \pageref{q:C3}\\
\ref{q:C4} & $J_2$ and $J_{22}$: meridians and parallels & list, Sep 2024 & \pageref{q:C4}\\
\ref{q:C5} & Three types of harmonics; $J_{32}$ & list, Sep 2024 & \pageref{q:C5}\\
\ref{q:C6} & Sun-synchronous orbits & list & \pageref{q:C6}\\
\ref{q:C7} & Third body: exact acceleration, effect on $\underline h$ & list, Jan 2025 & \pageref{q:C7}\\
\ref{q:C8} & Gravity-gradient dyadic & $\times$2, Feb 2025, Sep 2025 & \pageref{q:C8}\\
\ref{q:C9} & Drag: acceleration and effect on $a$, $e$ & Jul 2025, Mar 2025 & \pageref{q:C9}\\
\bottomrule
\end{tabular}
\end{center}

% ---------------------------------------------------------------------
\question{C1}{Solar radiation pressure: physical nature and perturbing acceleration}
% source: cap08.tex l.41-60 (intro), l.561-687 (SRP), l.714-726 (effect)
\begin{qbox}
Describe the physical nature of solar-radiation pressure; write the fundamental formula for the perturbing acceleration,
identifying and explaining each term.
\qmeta{$\times$3 in the list (2025); Jan 2023.}
\qnote{\qref{C2} is the same question plus the geometry of the Earth's shadow.}
\end{qbox}

\begin{outline}
\begin{enumerate}
\item Nature: electromagnetic radiation carries momentum; absorbed/reflected photons exert a pressure on the surface.
\item Intensity: $S=S_0(R_0/R)^2$; at 1\,AU $S_E=1367$\,W/m$^2$; pressure $P_{SR}=S_E/c=4.56\cdot10^{-6}$\,N/m$^2$.
\item Formula $\underline f^{(SR)}=-\nu\,\frac{P_{SR}A}{m}\,C_R\,\hat r_{SUN}$ (cannon-ball model) and meaning of $\nu$, $A$, $m$, $C_R\in[1,2]$, $\hat r_{SUN}$.
\item Characteristics: surface, independent of altitude, important at high orbits; effect: all elements, $\underline e$ orthogonal to the Sun direction.
\end{enumerate}
\end{outline}

\answer
\paragraph{Physical nature.} Solar radiation is electromagnetic radiation, and photons carry linear momentum ($p=E/c$). When
radiation hits the spacecraft surface, the absorbed photons transfer their momentum and the reflected ones transfer about
twice as much: the result is a small force on the illuminated surface, i.e.\ a \emph{pressure}. It is a \textbf{surface} perturbation
(it depends on the area exposed and on the optical properties of the surface) and, unlike all other perturbations, it does
\textbf{not depend on the altitude} of the orbit, only on the distance from the Sun.

\paragraph{Radiation pressure.} The radiated power intensity at the photosphere is $S_0=63.15\cdot10^6$\,W/m$^2$; radiation follows the
inverse-square law, so at a distance $R$ from the Sun centre
\begin{equation}
S=S_{0}\left(\frac{R_{0}}{R}\right)^{2},\qquad R_0=696\,000\ \text{km (photospheric radius)}.
\end{equation}
At the Earth's distance ($1\,\text{AU}=149.6\cdot10^6$\,km) $S_E=1367$\,W/m$^2$ (\emph{solar constant}). Dividing by the speed of light:
\begin{equation}
P_{SR}=\frac{S_{E}}{c}=4.56\cdot10^{-6}\ \frac{\text{N}}{\text{m}^{2}}.
\end{equation}

\paragraph{Perturbing acceleration (cannon-ball model).} The spacecraft is modelled as a sphere of radius $\sigma$:
\begin{equation}
\boxed{\;\underline{f}^{(SR)}=-\nu\,\frac{P_{SR}\,A}{m}\,C_{R}\,\hat{r}_{\text{SUN}}\;}
\end{equation}
\begin{itemize}
\item $\nu$ = \textbf{shadow function}: $\nu=1$ if the spacecraft is illuminated, $\nu=0$ if it is in the Earth's shadow (eclipse).
\item $P_{SR}$ = solar radiation pressure at the spacecraft distance from the Sun (4.56\,$\mu$N/m$^2$ at 1\,AU).
\item $A=\pi\sigma^2$ = \textbf{illuminated cross-section} (area projected on the plane normal to the Sun rays).
\item $m$ = spacecraft mass; $A/m$ (area-to-mass ratio) measures the sensitivity to SRP: large, light structures (solar sails,
balloons, large panels) are strongly affected.
\item $C_R$ = \textbf{radiation-pressure coefficient}, between 1 (totally absorbing surface) and 2 (totally reflecting surface). The
limits follow from conservation of linear momentum: a reflected photon reverses its momentum, so it transfers twice the
momentum of an absorbed one.
\item $\hat r_{SUN}$ = unit vector \textbf{from the spacecraft to the Sun}; the minus sign means the acceleration points away
from the Sun. Since $r\ll1$\,AU it can be approximated with the Earth-to-Sun direction:
$\underline r_{SUN}=1\,\text{AU}\,[c_{\theta_s}\ \ c_\epsilon s_{\theta_s}\ \ s_\epsilon s_{\theta_s}]\,[\hat c_1\ \hat c_2\ \hat c_3]^T$
($\epsilon$ = obliquity of the ecliptic, $\theta_s$ = ecliptic longitude of the Sun, increasing by $2\pi$ per sidereal year).
\end{itemize}
\begin{center}\includegraphics[width=0.62\linewidth]{srp}\end{center}

\paragraph{Effect on the orbit.} Projected on $(\hat r,\hat\theta,\hat h)$ the acceleration enters the Gauss equations. SRP perturbs
\textbf{all orbit elements}. Its relative importance grows with altitude (at 1000\,km $a_{RP}\sim10^{-9}a_0$; at GEO it is one of the main
perturbations together with the third body). For an initially circular orbit SRP adds velocity on one side of the orbit
(positive tangential acceleration, $P_1$) and removes it on the opposite side ($P_2$): the orbit becomes elliptic with an
\textbf{eccentricity vector orthogonal to $\hat r_{SUN}$}, which then precesses following the yearly motion of the Earth around the Sun.

\begin{keyf}
$S=S_0(R_0/R)^2$, \ $S_E=1367$\,W/m$^2$, \ $P_{SR}=S_E/c=4.56\cdot10^{-6}$\,N/m$^2$, \
$\underline f^{(SR)}=-\nu\dfrac{P_{SR}A}{m}C_R\hat r_{SUN}$, \ $1\le C_R\le2$, \ $A=\pi\sigma^2$.
\end{keyf}

\pitfalls
\begin{itemize}
\item $\hat r_{SUN}$ points from the spacecraft to the Sun: hence the minus sign (the push is away from the Sun).
\item $C_R=1$ absorbing, $C_R=2$ reflecting (not the reverse); justify with momentum conservation.
\item Do not forget the shadow function $\nu$: no SRP in eclipse.
\end{itemize}

\source{cap08.tex l.41--60 (classification), l.561--687 (SRP and Sun position), l.714--726 (effect on the elements). Formulario: not included (T only).}

% ---------------------------------------------------------------------
\question{C2}{Solar radiation pressure and the geometry of eclipses}
% source: cap08.tex l.561-726
\begin{qbox}
Outline the physical nature of solar radiation perturbation, provide the expression of the perturbing acceleration, and describe
the geometry of eclipsing due to Earth.
\qmeta{Jun 2025.}
\qnote{Difference from \qref{C1}: adds the shadow (eclipse) computation, i.e.\ how $\nu$ is evaluated.}
\end{qbox}

\begin{outline}
\begin{enumerate}
\item Nature (photon momentum, surface perturbation, independent of altitude); $S$, $S_E$, $P_{SR}=S_E/c$.
\item $\underline f^{(SR)}=-\nu\frac{P_{SR}A}{m}C_R\hat r_{SUN}$ and meaning of the terms.
\item Eclipse geometry: angle $\Theta$ between $\underline r$ and $\underline r_{SUN}$; $\theta_1=\arccos(R_E/r)$, $\theta_2=\arccos(R_E/r_{SUN})$; sketch.
\item Test: $\Theta>\theta_1+\theta_2$ shadow ($\nu=0$), otherwise illuminated ($\nu=1$); comments.
\end{enumerate}
\end{outline}

\answer
\paragraph{Physical nature and acceleration.} Photons carry momentum; when absorbed or reflected by the spacecraft surface they
exert a pressure. At distance $R$ from the Sun $S=S_0(R_0/R)^2$; at 1\,AU $S_E=1367$\,W/m$^2$ and the pressure is
$P_{SR}=S_E/c=4.56\cdot10^{-6}$\,N/m$^2$. With the cannon-ball model (sphere of radius $\sigma$):
\begin{equation}
\underline{f}^{(SR)}=-\nu\,\frac{P_{SR}\,A}{m}\,C_{R}\,\hat{r}_{\text{SUN}},
\end{equation}
$\nu$ shadow function (1 illuminated, 0 in shadow), $A=\pi\sigma^2$ illuminated area, $m$ mass, $C_R\in[1,2]$ radiation-pressure
coefficient (1 absorbing, 2 reflecting: a reflected photon transfers twice the momentum), $\hat r_{SUN}$ unit vector from the
spacecraft to the Sun ($\simeq$ Earth--Sun direction since $r\ll1$\,AU). It is a surface perturbation, independent of the
altitude, relevant for high orbits and for high $A/m$; it perturbs all elements and, on a circular orbit, creates an
eccentricity vector orthogonal to the Sun direction (see \qref{C1}).

\paragraph{Geometry of the eclipse (shadow computation).} The spacecraft is in shadow when the Earth (sphere of radius $R_E$) blocks
the line of sight between spacecraft and Sun. Let
\begin{equation}
\Theta=\arccos\!\left(\frac{\underline r\cdot\underline r_{SUN}}{|\underline r|\,|\underline r_{SUN}|}\right)
\end{equation}
be the angle, at the Earth's centre, between the spacecraft position $\underline r$ and the Sun position $\underline r_{SUN}$. Draw from the
spacecraft and from the Sun the lines tangent to the Earth: the radius to each tangent point is orthogonal to the tangent, so
\begin{equation}
\theta_{1}=\arccos\!\left(\frac{R_{E}}{|\underline r|}\right),\qquad
\theta_{2}=\arccos\!\left(\frac{R_{E}}{|\underline r_{SUN}|}\right)
\end{equation}
are the angles between $\underline r$ (resp.\ $\underline r_{SUN}$) and the radius to its tangent point. The segment spacecraft--Sun is tangent to the
Earth when $\Theta=\theta_1+\theta_2$. Hence
\begin{equation}
\boxed{\;\Theta>\theta_1+\theta_2\ \Rightarrow\ \text{SHADOW }(\nu=0),\qquad \Theta\le\theta_1+\theta_2\ \Rightarrow\ \text{ILLUMINATION }(\nu=1).\;}
\end{equation}
\begin{center}\includegraphics[width=0.9\linewidth]{eclipse}\\
{\footnotesize Not to scale ($r_{SUN}$ drawn small): in reality $\theta_2\simeq90^\circ$.}\end{center}

\paragraph{Comments.}
\begin{itemize}
\item Since $|\underline r_{SUN}|\gg R_E$, $\theta_2\simeq\pi/2$: the shadow is practically a cylinder of radius $R_E$ behind the Earth, and the
test becomes $\Theta>\pi/2+\theta_1$ \notinnotes.
\item The model is a sharp on/off shadow (no penumbra, spherical Earth, no atmospheric refraction) \notinnotes.
\item Eclipses make $\underline f^{(SR)}$ discontinuous along the orbit. Their duration is maximum when the Sun lies in the orbit plane:
for a circular orbit the shadow arc is $\pi-2\theta_1$, e.g.\ $r=7000$\,km: $\theta_1=24.3^\circ$, about 36\% of the period
(35\,min) \notinnotes. High orbits (GEO) are eclipsed only near the equinoxes.
\item Eclipses also matter for power (batteries) and thermal design.
\end{itemize}

\begin{fig2draw}
Earth circle with centre $O$; arrow $\underline r_{SUN}$ to the right, arrow $\underline r$ up-left; dashed tangent lines from the tip of each
arrow to the circle; mark $\theta_1$, $\theta_2$ at $O$ between each arrow and its tangent-point radius, and $\Theta$ between the arrows.
\end{fig2draw}

\begin{keyf}
$\underline f^{(SR)}=-\nu\dfrac{P_{SR}A}{m}C_R\hat r_{SUN}$; \ $\Theta=\arccos\dfrac{\underline r\cdot\underline r_{SUN}}{r\,r_{SUN}}$, \
$\theta_1=\arccos\dfrac{R_E}{r}$, \ $\theta_2=\arccos\dfrac{R_E}{r_{SUN}}$; \ shadow iff $\Theta>\theta_1+\theta_2$.
\end{keyf}

\pitfalls
\begin{itemize}
\item The angles $\theta_1$, $\theta_2$ are measured at the Earth's centre, between the position vectors and the radii to the tangent points.
\item Shadow for \emph{large} $\Theta$ (spacecraft behind the Earth), not small.
\end{itemize}

\source{cap08.tex l.561--687 (SRP), l.689--712 (shadow computation), l.714--726 (effect on the elements).}

% ---------------------------------------------------------------------
\question{C3}{$J_2$ perturbation: physical nature and averaged effects}
% source: cap08.tex l.728-836 (harmonics), l.838-949 (J2), l.951-1062 (averaging)
\begin{qbox}
Describe the physical nature of the $J_2$ perturbation, its two main averaged effects on orbital elements, and state whether it
depends on geographic latitude and/or longitude.
\qmeta{list by chapter (Ch.\,8).}
\qnote{Possibly the "question about $\Omega$" of Sep 2024 (regression of the node). Adjacent: \qref{C4}, \qref{C5}, \qref{C6}.}
\end{qbox}

\begin{outline}
\begin{enumerate}
\item Geopotential expansion; $J_2$ is the zonal harmonic of degree 2: Earth \textbf{oblateness} (equatorial bulge); dominant ($J_2=1.083\cdot10^{-3}$).
\item $U_{J2}=-\frac{\mu}{r}\left(\frac{R_E}{r}\right)^2J_2\frac{3s_\phi^2-1}{2}$: depends on latitude (and $r$), \textbf{not on longitude}; sign map ($\pm35.3^\circ$).
\item Perturbing acceleration (radial + northward) and averaging over one orbit: $\langle\dot a\rangle=\langle\dot e\rangle=\langle\dot i\rangle=0$.
\item Effect 1: precession of the orbit plane, $\langle\dot\Omega\rangle\propto-c_i$ (westward for direct orbits). Effect 2: rotation of the apsidal line, $\langle\dot\omega\rangle$; critical inclination $63.4^\circ$.
\item Applications: sun-synchronous orbits, frozen apsidal line (Molniya).
\end{enumerate}
\end{outline}

\answer
\paragraph{Physical nature.} The Earth is not a sphere with spherical mass distribution: its potential is expanded in harmonics
(Legendre polynomials $P_{lm}$),
\begin{equation}
U=\frac{\mu_E}{r}-\frac{\mu_E}{r}\sum_{l=2}^{\infty}\left(\frac{R_E}{r}\right)^{l}J_lP_{l0}(s_\phi)
+\frac{\mu_E}{r}\sum_{l=2}^{\infty}\sum_{m=1}^{l}\left(\frac{R_E}{r}\right)^{l}J_{lm}P_{lm}(s_\phi)\cos\!\left[m(\lambda_g-\lambda_{lm})\right].
\end{equation}
The first term is the Keplerian potential; $J_2$ (zonal harmonic of degree 2, $P_{20}=\frac{3s_\phi^2-1}{2}$) represents the
\textbf{oblateness} of the Earth, i.e.\ the equatorial bulge due to its rotation:
\begin{equation}
U_{J2}=-\frac{\mu_{E}}{r}\left(\frac{R_{E}}{r}\right)^{2}J_{2}\,\frac{3s_{\phi}^{2}-1}{2},\qquad J_2=1.083\cdot10^{-3}.
\end{equation}
$U_{J2}>0$ for $|\phi|<\arcsin(1/\sqrt3)=35.3^\circ$ (extra mass around the equator) and $U_{J2}<0$ for higher latitudes (less mass at
the poles). $J_2$ is about three orders of magnitude larger than any other harmonic, so it is the dominant perturbation in LEO and
MEO (at 1000\,km, $a_{J2}\simeq10^{-2}a_0$), decreasing rapidly with altitude ($\propto1/r^4$).

\paragraph{Dependence on latitude and longitude.} $U_{J2}$ depends only on $r$ and on the \textbf{latitude} $\phi$; it does \textbf{not depend on the
longitude} ($\partial U_{J2}/\partial\lambda_g=0$): the field is axially symmetric about the polar axis (zonal harmonic). Hence
the perturbing acceleration has only radial and northward components:
\begin{equation}
\underline f_{J2}=\frac{3\mu_E}{r^4}R_E^2J_2\left[\frac{3s_\phi^2-1}{2}\,\hat r-s_\phi c_\phi\,\hat N_L\right]
=\frac{3\mu_{E}}{r^{4}}R_{E}^{2}J_{2}\left[\frac{3s_{\theta_t}^{2}s_{i}^{2}-1}{2}\ \ \ -s_{i}^{2}s_{\theta_t}c_{\theta_t}\ \ \ -s_{i}c_{i}s_{\theta_t}\right]^{\!T}
\end{equation}
(components along $\hat r,\hat\theta,\hat h$). Through the orbit the satellite sweeps latitudes $|\phi|\le i$, so the effect depends on
the inclination.

\paragraph{Averaging.} The osculating elements show short-period oscillations, long-period oscillations and secular drifts. Averaging
the Gauss equations over one orbital period (changing variable from $t$ to $\theta_*$, $d/d\theta_*\simeq(r^2/h)\,d/dt$) removes the
short-period terms. Example for $\Omega$: $\langle\Omega'\rangle=\frac{1}{2\pi}\int_0^{2\pi}\frac{r^2}{h}\frac{rf_h s_{\theta_t}}{hs_i}d\theta_*=-\frac{3R_E^2J_2c_i}{2p^2}$, then
multiplied by the mean motion. The results are
\begin{equation}
\langle\dot a\rangle=\langle\dot e\rangle=\langle\dot i\rangle=0,\qquad
\langle\dot{\Omega}\rangle=-\frac{3R_{E}^{2}J_{2}\sqrt{\mu_{E}}}{2a^{7/2}(1-e^{2})^{2}}\,c_{i},\qquad
\langle\dot{\omega}\rangle=-\frac{3R_{E}^{2}J_{2}\sqrt{\mu_{E}}}{2a^{7/2}(1-e^{2})^{2}}\left(\frac52s_{i}^{2}-2\right).
\end{equation}
Size and shape of the orbit and its inclination are unchanged on average (conservative, axisymmetric field).

\paragraph{The two main effects.}
\begin{enumerate}
\item \textbf{Precession of the orbital plane (regression of the nodes).} $\langle\dot\Omega\rangle\neq0$ with $i$ constant: the angular momentum
vector rotates about $\hat c_3$ (the Earth's axis) with period $T_{prec}=2\pi/|\langle\dot\Omega\rangle|$. The precession is clockwise (westward,
$\langle\dot\Omega\rangle<0$) for direct orbits ($i<\pi/2$), counter-clockwise for retrograde ones, and zero for polar orbits. Physically, the
equatorial bulge exerts a torque that tries to pull the orbit plane towards the equator; like a gyroscope, the plane precesses
instead of changing $i$ \notinnotes. Faster for low, near-equatorial orbits (several degrees per day in LEO).
\item \textbf{Rotation of the apsidal line.} $\langle\dot\omega\rangle\neq0$: the periapsis direction rotates within the orbit plane (forward for
$i<63.4^\circ$, backward between $63.4^\circ$ and $116.6^\circ$). It vanishes at the \textbf{critical inclinations}
\begin{equation}
5s_i^2-4=0\ \Rightarrow\ i=\arcsin\frac{2}{\sqrt5}\simeq63.4^\circ\ \text{or}\ 116.6^\circ,
\end{equation}
where the apsidal line is "frozen" in space (used by highly elliptic Molniya orbits to keep the apogee over high latitudes \notinnotes).
\end{enumerate}
\begin{center}\includegraphics[width=0.85\linewidth]{j2_effects}\end{center}
Application: the nodal precession can be tuned to $360^\circ$/year to obtain \textbf{sun-synchronous orbits} (\qref{C6}).

\begin{keyf}
$U_{J2}=-\dfrac{\mu_E}{r}\left(\dfrac{R_E}{r}\right)^2J_2\dfrac{3s_\phi^2-1}{2}$; \ $\langle\dot a\rangle=\langle\dot e\rangle=\langle\dot i\rangle=0$; \
$\langle\dot\Omega\rangle=-\dfrac{3R_E^2J_2\sqrt{\mu_E}}{2a^{7/2}(1-e^2)^2}c_i$; \
$\langle\dot\omega\rangle=-\dfrac{3R_E^2J_2\sqrt{\mu_E}}{2a^{7/2}(1-e^2)^2}\left(\frac52s_i^2-2\right)$; \ $i_{crit}=63.4^\circ,\,116.6^\circ$.
\end{keyf}

\pitfalls
\begin{itemize}
\item $J_2$ depends on latitude only (zonal), not on longitude: say it explicitly with $\partial U_{J2}/\partial\lambda_g=0$.
\item Sign of $\langle\dot\Omega\rangle$: negative (westward) for direct orbits.
\item The secular effects are on $\Omega$ and $\omega$ (and on the mean motion), not on $a$, $e$, $i$.
\end{itemize}

\source{cap08.tex l.728--836 (geopotential harmonics), l.838--868 ($J_2$), l.870--949 (perturbing acceleration; errata \#12--13),
l.951--1062 (averaging, precession, critical inclinations). Formulario: \fml{8-J2V}, \fml{8-J2A}, \fml{8-PRC}.}

% ---------------------------------------------------------------------
\question{C4}{Physical nature of $J_2$ and $J_{22}$; meridians and parallels}
% source: cap08.tex l.728-868
\begin{qbox}
State the physical nature of $J_2$ and $J_{22}$ and how many meridians/parallels each corresponds to.
\qmeta{list by chapter (Ch.\,8); Sep 2024 (together with \qref{C5}).}
\end{qbox}

\begin{outline}
\begin{enumerate}
\item Geopotential expansion; $J_{lm}$ terms; rule: $J_{lm}$ has $l-m$ parallels and $m$ meridians as nodal (equipotential) lines.
\item $J_2=J_{20}$: zonal, Earth oblateness (equatorial bulge); $P_{20}=0$ at $\phi=\pm35.3^\circ$: 2 parallels, 0 meridians; depends on latitude only.
\item $J_{22}$: sectoral, ellipticity of the equator; zero at $\lambda_g=30.1^\circ,120.1^\circ,210.1^\circ,300.1^\circ$: 2 meridians, 0 parallels; depends on longitude.
\item Magnitudes and effects ($J_2$ dominant in LEO/MEO; $J_{22}$ dominant at GEO by resonance). Sketch.
\end{enumerate}
\end{outline}

\answer
\paragraph{Geopotential.} The Earth's gravitational potential is written as a series of harmonics:
\begin{equation}
U=\frac{\mu_E}{r}-\frac{\mu_E}{r}\sum_{l\ge2}\left(\frac{R_E}{r}\right)^{l}J_lP_{l0}(s_\phi)
+\frac{\mu_E}{r}\sum_{l\ge2}\sum_{m=1}^{l}\left(\frac{R_E}{r}\right)^{l}J_{lm}P_{lm}(s_\phi)\cos\!\left[m(\lambda_g-\lambda_{lm})\right],
\end{equation}
\begin{equation}
P_{lm}(x)=\frac{1}{2^{l}l!}\left(1-x^{2}\right)^{m/2}\frac{d^{\,l+m}}{dx^{\,l+m}}\left[(x^{2}-1)^{l}\right],\qquad x=s_\phi .
\end{equation}
Each harmonic is a correction to the potential of a spherical Earth; its nodal lines (where the term vanishes) are parallels and/or
meridians. General rule: \textbf{$J_{lm}$ has $l-m$ parallels and $m$ meridians} as equipotential (zero) lines.

\paragraph{$J_2$ ($l=2$, $m=0$): zonal harmonic -- Earth oblateness.} Since $P_{20}=\frac{3s_\phi^2-1}{2}$,
\[U_{J2}=-\frac{\mu_E}{r}\left(\frac{R_E}{r}\right)^2J_2\,\frac{3s_\phi^2-1}{2}\]
depends only on latitude. It vanishes where
$s_\phi=\pm1/\sqrt3$, i.e.\ at $\phi=\pm35.3^\circ$: \textbf{2 parallels and 0 meridians} ($l-m=2$, $m=0$). It is positive in the
equatorial band and negative at high latitudes: it models the \textbf{equatorial bulge} (extra mass at the equator, less at the poles)
produced by the Earth's rotation. $J_2=1.083\cdot10^{-3}$, about $10^3$ times larger than any other coefficient: it dominates the
perturbations of LEO and MEO satellites (nodal regression, apsidal rotation, \qref{C3}).

\paragraph{$J_{22}$ ($l=m=2$): sectoral harmonic -- ellipticity of the equator.} $P_{22}=3(1-s_\phi^2)=3c_\phi^2$ never vanishes except at the
poles, so the sign of the term is set by $\cos[2(\lambda_g-\lambda_{22})]$: it depends on longitude and vanishes at
\begin{equation}
\lambda_g=30.1^\circ,\ 120.1^\circ,\ 210.1^\circ,\ 300.1^\circ\quad(90^\circ\text{ apart}),
\end{equation}
i.e.\ on \textbf{2 meridians (full great circles) and 0 parallels} ($l-m=0$, $m=2$), dividing the Earth into 4 sectors of alternating sign
("orange slices"). It is related to the modest \textbf{eccentricity (ellipticity) of the Earth's equator}: the equatorial section is slightly
elliptic, not circular. $J_{22}=1.816\cdot10^{-6}$ ($\lambda_{22}=-0.261$\,rad). Small, but it \textbf{dominates for geostationary orbits}: a GEO
satellite is fixed with respect to the Earth, so the longitude-dependent force always acts in the same direction (resonance) and
makes the satellite drift in longitude (station keeping needed).
\begin{center}\includegraphics[width=0.9\linewidth]{harmonics}\\
{\footnotesize Sign maps of the three types of harmonics (hatched = negative), computed from $P_{lm}(s_\phi)\cos m\lambda$.}\end{center}

\begin{center}
\small
\begin{tabular}{@{}llllll@{}}
\toprule
 & type & physical meaning & depends on & parallels & meridians\\
\midrule
$J_2$ & zonal ($m=0$) & oblateness (equatorial bulge) & latitude & 2 ($\pm35.3^\circ$) & 0\\
$J_{22}$ & sectoral ($l=m$) & elliptic equator & longitude & 0 & 2 (4 sectors)\\
\bottomrule
\end{tabular}
\end{center}

\begin{keyf}
$J_{lm}$: $l-m$ parallels, $m$ meridians; \ $P_{20}=\frac{3s_\phi^2-1}{2}=0\Rightarrow\phi=\pm35.3^\circ$; \ $J_{22}$ zero at $\lambda_g=30.1^\circ+k\,90^\circ$; \
$J_2=1.083\cdot10^{-3}$, $J_{22}=1.816\cdot10^{-6}$.
\end{keyf}

\pitfalls
\begin{itemize}
\item "2 meridians" for $J_{22}$ means 2 great circles through the poles, i.e.\ 4 longitudes where the term vanishes.
\item $J_2$ = oblateness (polar flattening), $J_{22}$ = ellipticity of the \emph{equator}: different physical features.
\end{itemize}

\source{cap08.tex l.62--85 (perturbations on different orbits, $J_{22}$ at GEO), l.728--836 (harmonics, nodal lines, coefficients), l.838--868 ($J_2$).}

% ---------------------------------------------------------------------
\question{C5}{The three types of geopotential harmonics; the $J_{32}$ harmonic}
% source: cap08.tex l.728-836
\begin{qbox}
Describe the three types of geopotential harmonics; indicate whether they depend on latitude, longitude, or both, and give the
number and type of equipotential lines for the $J_{32}$ harmonic.
\qmeta{list by chapter (Ch.\,8); Sep 2024 (together with \qref{C4}: "also discuss the physical nature of $J_2$ and $J_{22}$").}
\qnote{In Sep 2024 this question was asked together with C4: write C5 and then the two paragraphs on $J_2$ and $J_{22}$ of C4.}
\end{qbox}

\begin{outline}
\begin{enumerate}
\item Geopotential series with $J_l$ and $J_{lm}$, $P_{lm}(s_\phi)$, $\cos m(\lambda_g-\lambda_{lm})$; first term = spherical Earth.
\item Zonal ($m=0$): latitude only, $l$ parallels (bands); e.g.\ $J_2$.
\item Sectoral ($l=m$): longitude only (sign), $m$ meridians (sectors); e.g.\ $J_{22}$.
\item Tesseral ($0<m<l$): latitude and longitude, $l-m$ parallels and $m$ meridians (tiles).
\item $J_{32}$: tesseral, 1 parallel (equator) + 2 meridians, 8 tiles; sketch.
\end{enumerate}
\end{outline}

\answer
\paragraph{Expansion of the geopotential.} The potential of the Earth, $U=G\int_{Earth}dm/r$, is written in terms of Legendre polynomials
(latitude $\phi$, geographical longitude $\lambda_g$):
\begin{equation}
U=\underbrace{\frac{\mu_E}{r}}_{U_K}
-\frac{\mu_E}{r}\sum_{l=2}^{\infty}\left(\frac{R_E}{r}\right)^{l}J_lP_{l0}(s_\phi)
+\frac{\mu_E}{r}\sum_{l=2}^{\infty}\sum_{m=1}^{l}\left(\frac{R_E}{r}\right)^{l}J_{lm}P_{lm}(s_\phi)\cos\!\left[m(\lambda_g-\lambda_{lm})\right].
\end{equation}
$U_K=\mu_E/r$ is the potential of a body with spherical symmetry (Keplerian motion); the other terms are the harmonics, with
coefficients $J_l$, $J_{lm}$, $\lambda_{lm}$ estimated from satellite tracking (WGS-84, EGM-96, JGM-2, EGM2008). The perturbing
acceleration is $\underline g=\nabla U$.

\paragraph{The three types.}
\begin{enumerate}[label=(\alph*)]
\item \textbf{Zonal harmonics} ($m=0$, coefficients $J_l$): $P_{l0}(s_\phi)$ \textbf{depends only on latitude}. They vanish at $l$ values of latitude:
the nodal lines are \textbf{$l$ parallels} and the Earth is divided in latitude bands (zones). Example: $J_2$ vanishes at $\pm35.3^\circ$
(oblateness).
\item \textbf{Sectoral harmonics} ($l=m$, coefficients $J_{ll}$): $P_{ll}\propto c_\phi^{\,l}$ does not vanish except at the poles, so the sign of the
term changes only with longitude: they \textbf{depend (in sign) only on longitude}. Nodal lines: \textbf{$m$ meridians}, which split the Earth
into $2m$ sectors. Example: $J_{22}$, zero at $\lambda_g=30.1^\circ,120.1^\circ,210.1^\circ,300.1^\circ$ (ellipticity of the equator).
\item \textbf{Tesseral harmonics} ($0<m<l$, coefficients $J_{lm}$): \textbf{depend on both latitude and longitude}; they vanish at given latitudes
\emph{and} longitudes, so the nodal lines (\textbf{$l-m$ parallels and $m$ meridians}) divide the Earth into tiles ("tesserae").
\end{enumerate}
General rule: \textbf{$J_{lm}$ has $l-m$ parallels and $m$ meridians as equipotential lines} (zonal: $m=0$; sectoral: $l-m=0$).

\paragraph{The $J_{32}$ harmonic.} $l=3$, $m=2$: \textbf{tesseral}. From the definition,
$P_{32}(x)=\frac{1}{48}(1-x^2)\frac{d^5}{dx^5}(x^2-1)^3=15x(1-x^2)$, so the term is proportional to
\begin{equation}
15\,s_\phi c_\phi^2\cos\!\left[2(\lambda_g-\lambda_{32})\right],
\end{equation}
which vanishes on
\begin{itemize}
\item $l-m=1$ \textbf{parallel}: the \textbf{equator} ($s_\phi=0$);
\item $m=2$ \textbf{meridians}: $\lambda_g=\lambda_{32}\pm45^\circ$ and $\lambda_{32}\pm135^\circ$ (two great circles through the poles).
\end{itemize}
The equipotential lines divide the Earth into $2\times4=8$ tiles of alternating sign. $J_{32}=3.774\cdot10^{-6}$,
$\lambda_{32}=-0.300$\,rad.
\begin{center}\includegraphics[width=0.9\linewidth]{harmonics}\end{center}

\begin{center}
\small
\begin{tabular}{@{}lllll@{}}
\toprule
type & indices & depends on & nodal lines & example\\
\midrule
zonal & $m=0$ & latitude & $l$ parallels & $J_2$: $\pm35.3^\circ$\\
sectoral & $m=l$ & longitude & $m$ meridians & $J_{22}$: 2 meridians\\
tesseral & $0<m<l$ & both & $l-m$ parallels, $m$ meridians & $J_{32}$: equator + 2 meridians\\
\bottomrule
\end{tabular}
\end{center}

\begin{fig2draw}
Three circles (globes). Zonal: two horizontal lines (parallels) symmetric about the equator, polar caps hatched. Sectoral: two
"meridian" curves through the poles, alternate sectors hatched. Tesseral $J_{32}$: the equator plus two meridians, checkerboard hatching.
\end{fig2draw}

\begin{keyf}
$U=\dfrac{\mu_E}{r}-\dfrac{\mu_E}{r}\displaystyle\sum_l\left(\frac{R_E}{r}\right)^lJ_lP_{l0}+\dfrac{\mu_E}{r}\sum_{l,m}\left(\frac{R_E}{r}\right)^lJ_{lm}P_{lm}\cos m(\lambda_g-\lambda_{lm})$; \
$J_{lm}$: $l-m$ parallels, $m$ meridians; \ $J_{32}$: 1 parallel (equator) + 2 meridians.
\end{keyf}

\pitfalls
\begin{itemize}
\item Count meridians as great circles (each meridian gives two zero longitudes $180^\circ$ apart).
\item "Sectoral depends only on longitude" refers to the sign / nodal lines; the magnitude still decreases towards the poles.
\end{itemize}

\source{cap08.tex l.728--836 (harmonics, $P_{lm}$, types, nodal-line rule, EGM2008 coefficients). $P_{32}$ computed from the formula of the notes.}

% ---------------------------------------------------------------------
\question{C6}{Sun-synchronous orbits}
% source: cap08.tex l.1029-1099
\begin{qbox}
Define a Sun-synchronous orbit about Earth and name the natural perturbation used to achieve Sun-synchronism. Specify whether
Sun-synchronous orbits are prograde or retrograde and justify your answer.
\qmeta{list by chapter (Ch.\,8).}
\end{qbox}

\begin{outline}
\begin{enumerate}
\item Definition: the orbit plane rotates with the same angular rate as the Earth around the Sun ($360^\circ$/year) $\Rightarrow$ constant geometry w.r.t.\ the Sun, near-identical lighting.
\item Perturbation: $J_2$ (oblateness) $\Rightarrow$ nodal precession $\langle\dot\Omega\rangle=-\frac{3R_E^2J_2\sqrt{\mu}}{2a^{7/2}(1-e^2)^2}c_i$.
\item Condition $\langle\dot\Omega\rangle=2\pi/T_{yr}$; formula for $i$.
\item Justification: precession must be counter-clockwise (eastward, like the Earth's motion around the Sun) $\Rightarrow$ $c_i<0$ $\Rightarrow$ \textbf{retrograde}, $i>90^\circ$ (about $97$--$99^\circ$ in LEO).
\end{enumerate}
\end{outline}

\answer
\paragraph{Definition.} A sun-synchronous orbit is an orbit whose plane \textbf{rotates with the same angular rate as the Earth around the
Sun}, i.e.\ one revolution per year ($\simeq0.9856^\circ$/day), in the same sense. Then the angle between the orbit plane and the Sun
direction (precisely: between the projection $\hat h_p$ of $\underline h$ on the ecliptic and the Sun direction) stays constant all year: the
satellite crosses each latitude at the same local solar time and sees \textbf{near-identical lighting conditions} (Earth observation,
remote sensing, constant illumination of solar panels, e.g.\ dawn--dusk orbits).
\begin{center}\includegraphics[width=0.85\linewidth]{sso}\end{center}

\paragraph{Perturbation used.} The \textbf{$J_2$ harmonic (Earth oblateness)}. Averaged over one orbit it produces a secular precession
of the node,
\begin{equation}
\langle\dot{\Omega}\rangle=-\frac{3R_{E}^{2}J_{2}\sqrt{\mu_{E}}}{2a^{7/2}(1-e^{2})^{2}}\,c_{i},\qquad
T_{prec}=\frac{2\pi}{|\langle\dot\Omega\rangle|}=\frac{4\pi\,a^{7/2}(1-e^2)^2}{3R_E^2J_2\sqrt{\mu_E}\,|c_i|},
\end{equation}
whose magnitude can be tuned with $a$, $e$, $i$. Sun-synchronism requires $T_{prec}=1$ year \emph{and} the correct sense of rotation:
\begin{equation}
\langle\dot{\Omega}\rangle=+\frac{2\pi}{T_{yr}}\ \Longrightarrow\
i=\arccos\!\left[-\frac{2\pi}{T_{yr}}\,\frac{2\,a^{7/2}(1-e^{2})^{2}}{3R_{E}^{2}J_{2}\sqrt{\mu_{E}}}\right].
\end{equation}

\paragraph{Prograde or retrograde?} \textbf{Retrograde.} The Earth moves around the Sun counter-clockwise as seen from the north
ecliptic pole, so the orbit plane must precess counter-clockwise (eastward): $\langle\dot\Omega\rangle>0$. Since the factor multiplying
$c_i$ is negative, this requires $c_i<0$, i.e.\ $i>\pi/2$. For direct orbits ($i<90^\circ$) $J_2$ makes the node regress westward,
the wrong way. Typical values for circular orbits: $i\simeq97.4^\circ$ at 500\,km, $98.2^\circ$ at 700\,km, $98.6^\circ$ at 800\,km
(near-polar, slightly retrograde); the inclination grows with altitude, and above $a\simeq12\,350$\,km ($c_i=-1$) no sun-synchronous
circular orbit exists \notinnotes (values computed from the formula above). The notes report the plot of altitude vs inclination for $e=0$.

\begin{keyf}
$\langle\dot\Omega\rangle=-\dfrac{3R_E^2J_2\sqrt{\mu_E}}{2a^{7/2}(1-e^2)^2}c_i=+\dfrac{2\pi}{T_{yr}}$ \ $\Rightarrow$ \
$c_i=-\dfrac{2\pi}{T_{yr}}\dfrac{2a^{7/2}(1-e^2)^2}{3R_E^2J_2\sqrt{\mu_E}}<0$ \ $\Rightarrow$ retrograde, $i\approx98^\circ$ in LEO.
\end{keyf}

\pitfalls
\begin{itemize}
\item The sense of precession matters (not only the period): this is what forces $i>90^\circ$.
\item Sun-synchronous is not "geosynchronous" and not "polar": it is slightly retrograde.
\item The perturbation is $J_2$ (oblateness), not the Sun's gravity.
\end{itemize}

\source{cap08.tex l.1029--1062 (precession of the plane), l.1064--1099 (sun-synchronous orbits). Formulario: \fml{8-SSO}.}

% ---------------------------------------------------------------------
\question{C7}{Third-body perturbation: exact acceleration and effect on the angular momentum}
% source: cap08.tex l.1101-1205 (exact), l.1211-1329 (gravity gradient), l.1331-1453 (average effect)
\begin{qbox}
Derive the exact third-body perturbing acceleration and describe how a third-body perturbation affects the angular momentum of
a spacecraft in a circular Earth orbit.
\qmeta{list by chapter ("2025's Q"); Jan 2025 ("physical nature of the third-body perturbation and derivation of the exact acceleration formula").}
\qnote{If only the nature and the exact formula are asked (Jan 2025): write points 1--3 of the outline. Adjacent: \qref{C8}.}
\end{qbox}

\begin{outline}
\begin{enumerate}
\item Physical nature: the third body (Sun, Moon) attracts both the spacecraft and the Earth; the perturbation is the \emph{difference} (tidal).
\item Inertial equations for the spacecraft and for the Earth; subtraction $\Rightarrow$ relative motion = Keplerian + $\underline a_{3B}$.
\item Exact formula $\underline a_{3B}=\mu_2\left[\frac{\underline r_{21}}{r_{21}^3}-\frac{\underline r_{2S}}{r_{2S}^3}\right]=-\mu_2\left\{\frac{\underline r_{12}}{r_{12}^3}+\frac{\underline r_{1S}-\underline r_{12}}{|\underline r_{1S}-\underline r_{12}|^3}\right\}$.
\item Effect on $\underline h$ (circular orbit): gravity gradient, $\dot{\underline h}=\underline r\times\underline f$, double averaging $\Rightarrow$ $\langle\langle\dot{\underline h}\rangle\rangle=\underline\omega_h\times\underline h$: precession about $\hat h_{3B}$, $|\underline h|$ constant.
\end{enumerate}
\end{outline}

\answer
\paragraph{Physical nature.} A spacecraft orbiting the Earth (body 1) is also attracted by other celestial bodies, mainly the Moon
and the Sun (body 2). The Earth itself is accelerated by body 2. What perturbs the motion \emph{relative to the Earth} is the
\textbf{difference} between the attraction of body 2 on the spacecraft and on the Earth (a tidal, differential effect). It is a
conservative, gravitational perturbation; it is small in LEO ($\simeq10^{-7}a_0$ at 1000\,km) and becomes a main perturbation
at high altitude (GEO and above).
\begin{center}\includegraphics[width=0.48\linewidth]{thirdbody}\end{center}

\paragraph{Derivation of the exact acceleration.} Notation: $\underline r_S$, $\underline r_1$, $\underline r_2$ inertial positions of spacecraft,
Earth and third body; $\underline r_{jS}=\underline r_S-\underline r_j$; $\underline r_{21}=\underline r_1-\underline r_2=-\underline r_{12}$;
$\mu_j$ gravitational parameters. In an inertial frame:
\begin{equation}
\frac{d^{2}\underline{r}_{S}}{dt^{2}}=-\frac{\mu_{1}}{r_{1S}^{3}}\underline{r}_{1S}-\frac{\mu_{2}}{r_{2S}^{3}}\underline{r}_{2S},
\qquad
\frac{d^{2}\underline{r}_{1}}{dt^{2}}=-\frac{\mu_{2}}{r_{21}^{3}}\underline{r}_{21}
\end{equation}
(the spacecraft exerts negligible gravity on the Earth). Subtracting, the motion relative to the Earth, $\underline r_{1S}=\underline r_S-\underline r_1$, obeys
\begin{equation}
\frac{d^{2}\underline{r}_{1S}}{dt^{2}}=-\frac{\mu_{1}}{r_{1S}^{3}}\underline{r}_{1S}
+\underbrace{\left[\frac{\mu_{2}}{r_{21}^{3}}\underline{r}_{21}-\frac{\mu_{2}}{r_{2S}^{3}}\underline{r}_{2S}\right]}_{\underline a_{3B}} .
\end{equation}
With $\underline r_{2S}=\underline r_{21}+\underline r_{1S}=\underline r_{1S}-\underline r_{12}$:
\begin{equation}
\boxed{\;\underline{a}_{3B}=-\mu_{2}\left\{\frac{\underline{r}_{12}}{r_{12}^{3}}
+\frac{\underline{r}_{1S}-\underline{r}_{12}}{\left[(\underline{r}_{1S}-\underline{r}_{12})\cdot(\underline{r}_{1S}-\underline{r}_{12})\right]^{3/2}}\right\}\;}
\end{equation}
The first term is the (indirect) acceleration of the Earth towards body 2, the second the (direct) attraction of body 2 on the
spacecraft. The formula is exact; to use it in the Gauss equations, $\underline r_{12}$ (Moon or Sun position w.r.t.\ the Earth) and
$\underline r_{1S}=\frac{p}{1+ec_{\theta_*}}\hat r$ are projected on $(\hat r,\hat\theta,\hat h)$.

\paragraph{Effect on the angular momentum of a circular orbit.}
\begin{enumerate}
\item Since $r=r_{1S}\ll r_{12}$, the acceleration is approximated by the \textbf{gravity gradient} (\qref{C8}):
$\underline f_{3B}\simeq\frac{\mu_2}{r_{12}^5}\left[3\underline r_{12}\underline r_{12}^T-r_{12}^2\underrightarrow{1}\right]\cdot\underline r$.
\item $\dot{\underline h}=\underline r\times\underline f_{3B}$; the term in $\underrightarrow{1}$ gives $\underline r\times\underline r=\underline 0$:
\begin{equation}
\frac{d\underline h}{dt}=\frac{3\mu_2}{r_{12}^5}(\underline r\times\underline r_{12})(\underline r_{12}\cdot\underline r)=-\frac{3\mu_2}{r_{12}^5}\,\underline r_{12}\times\left(\underline r\,\underline r^T\right)\underline r_{12}.
\end{equation}
\item \textbf{First averaging} over one spacecraft period ($r_{12}$ frozen, $\hat r=\hat Nc_{\theta_t}+\hat Ms_{\theta_t}$ in the orbit plane, $r=R$):
$\langle c^2_{\theta_t}\rangle=\langle s^2_{\theta_t}\rangle=\frac12$, $\langle c_{\theta_t}s_{\theta_t}\rangle=0$ $\Rightarrow$
$\langle\underline r\,\underline r^T\rangle=R^2\left[\frac12\underrightarrow{1}-\frac12\hat h\hat h^T\right]$, hence
$\left\langle\dot{\underline h}\right\rangle=\frac{3\mu_2R^2}{2r_{12}^5}(\underline r_{12}\times\hat h)(\hat h\cdot\underline r_{12})$.
\item \textbf{Second averaging} over one period of the perturbing body (circular relative orbit with normal $\hat h_{3B}$):
$\langle\underline r_{12}\underline r_{12}^T\rangle=r_{12}^2\left[\frac12\underrightarrow{1}-\frac12\hat h_{3B}\hat h_{3B}^T\right]$, which gives
\begin{equation}
\left\langle\!\left\langle\frac{d\underline{h}}{dt}\right\rangle\!\right\rangle=\underline\omega_h\times\underline h,\qquad
\underline{\omega}_{h}=-\frac{3}{4}\frac{\mu_{2}R^{2}}{r_{12}^{3}\,h}\left(\hat{h}\cdot\hat{h}_{3B}\right)\hat{h}_{3B},\qquad h=\sqrt{\mu_1R}.
\end{equation}
\end{enumerate}
\textbf{Result}: on average $\underline h$ only \textbf{rotates} about $\hat h_{3B}$ (normal to the orbit of the third body: ecliptic pole for the Sun),
\textbf{its magnitude does not change} (orbit size unchanged, $\dot{\underline h}\perp\underline h$), while the orbit plane precesses with
period $T_{prec}=2\pi/|\underline\omega_h|$; clockwise if $\hat h\cdot\hat h_{3B}>0$, counter-clockwise if $<0$. The inclination with respect to
the third body's orbit plane stays constant on average.
Example: the Moon, regarded as an Earth satellite perturbed by the Sun ($\hat h_{3B}$ = ecliptic pole): its node regresses with
$T_{prec}\simeq18$ years.

\begin{keyf}
$\underline a_{3B}=\mu_2\left[\dfrac{\underline r_{21}}{r_{21}^3}-\dfrac{\underline r_{2S}}{r_{2S}^3}\right]$, $\underline r_{2S}=\underline r_{21}+\underline r_{1S}$; \
$\dot{\underline h}=\underline r\times\underline f$; \ $\langle\langle\dot{\underline h}\rangle\rangle=\underline\omega_h\times\underline h$, \
$\underline\omega_h=-\dfrac34\dfrac{\mu_2R^2}{r_{12}^3h}(\hat h\cdot\hat h_{3B})\hat h_{3B}$.
\end{keyf}

\pitfalls
\begin{itemize}
\item The perturbation is the \emph{difference} of the two attractions (spacecraft minus Earth), not the attraction on the spacecraft alone.
\item Write the inertial equations for both bodies and subtract; keep track of $\underline r_{21}=-\underline r_{12}$.
\item The average effect is a precession of $\underline h$ at constant magnitude.
\end{itemize}

\source{cap08.tex l.1101--1205 (exact acceleration), l.1211--1329 (gravity gradient), l.1331--1453 (average effect on circular orbits;
errata \#14--15), l.1455--1470 (Moon). Formulario: \fml{8-3BA}, \fml{8-MOO}.}

% ---------------------------------------------------------------------
\question{C8}{Gravity-gradient dyadic from the third-body acceleration}
% source: cap08.tex l.1211-1329
\begin{qbox}
Prove the gravity-gradient dyadic starting from the third-body perturbing-force formula
\[
\underline f_{3B}=-\frac{\mu_S}{r_{2S}^3}\left(\underline r_{21}+\underline r_{1S}\right)+\frac{\mu_S}{r_{21}^3}\,\underline r_{21},
\]
while specifying the fundamental condition needed to obtain it.
\qmeta{$\times$2 in the list (2025); Feb 2025; Sep 2025.}
\qnote{Here body 2 is the Sun ($\mu_2=\mu_S$). Adjacent: \qref{C7}.}
\end{qbox}

\begin{outline}
\begin{enumerate}
\item Write $\underline f_{3B}=\underline g(\underline r_{21}+\underline r_{1S})-\underline g(\underline r_{21})$ with $\underline g(\underline x)=-\mu_2\underline x/x^3$.
\item Fundamental condition: $|\underline r_{1S}|\ll|\underline r_{21}|$ (spacecraft much closer to the Earth than the third body).
\item First-order Taylor expansion about $\underline r_{21}$: $\underline f_{3B}\simeq\left.\frac{\partial\underline g}{\partial\underline x}\right|_{\underline r_{21}}\cdot\underline r_{1S}$ (a dyad).
\item Components: $\partial_j(-\mu_2x_i/x^3)=-\mu_2\delta_{ij}/x^3+3\mu_2x_ix_j/x^5$ (show (1,1) and (1,2)).
\item Compact form $\frac{\mu_2}{r_{12}^5}\left[3\underline r_{12}\underline r_{12}-r_{12}^2\underrightarrow{1}\right]$; final $\underline f_{3B}$; interpretation (tidal stretching).
\end{enumerate}
\end{outline}

\answer
\paragraph{Starting point.} With $\underline r_{2S}=\underline r_{21}+\underline r_{1S}$ the third-body acceleration is the difference of the
gravitational field of body 2,
\begin{equation}
\underline g(\underline x):=-\frac{\mu_2}{x^3}\underline x,\qquad
\underline f_{3B}=-\frac{\mu_2}{r_{2S}^3}\underline r_{2S}-\left(-\frac{\mu_2}{r_{21}^3}\underline r_{21}\right)
=\underline g(\underline r_{21}+\underline r_{1S})-\underline g(\underline r_{21}),
\end{equation}
evaluated at the spacecraft ($\underline x=\underline r_{2S}$) and at the Earth ($\underline x=\underline r_{21}$).

\paragraph{Fundamental condition.}
\begin{equation}
\boxed{\;|\underline r_{1S}|\ll|\underline r_{21}|\;}
\end{equation}
the distance of the spacecraft from the Earth is much smaller than the distance of the third body from the Earth (usually satisfied:
LEO--GEO radii $\ll$ 384\,400\,km or 1\,AU). Then $\underline r_{1S}$ is a small displacement from $\underline r_{21}$ and a first-order Taylor
expansion is accurate:
\begin{equation}
\underline f_{3B}\simeq\left.\frac{\partial}{\partial\underline r_{2j}}\left[-\frac{\mu_2}{r_{2j}^3}\underline r_{2j}\right]\right|_{\underline r_{21}}\!\cdot\underline r_{1S}.
\end{equation}

\paragraph{The gravity-gradient dyad.} The gradient of a vector field with respect to a position vector is a dyad (second-order tensor),
which acting on a vector gives a vector. In an orthonormal basis $\{\hat c_1,\hat c_2,\hat c_3\}$, with $\underline r_{2j}=x\hat c_1+y\hat c_2+z\hat c_3$:
\begin{equation}
\underrightarrow{D}=[\hat c_1\ \hat c_2\ \hat c_3]
\begin{bmatrix}
\partial_x\!\left(-\frac{\mu_2x}{r_{2j}^3}\right)&\partial_y\!\left(-\frac{\mu_2x}{r_{2j}^3}\right)&\partial_z\!\left(-\frac{\mu_2x}{r_{2j}^3}\right)\\[3pt]
\partial_x\!\left(-\frac{\mu_2y}{r_{2j}^3}\right)&\partial_y\!\left(-\frac{\mu_2y}{r_{2j}^3}\right)&\partial_z\!\left(-\frac{\mu_2y}{r_{2j}^3}\right)\\[3pt]
\partial_x\!\left(-\frac{\mu_2z}{r_{2j}^3}\right)&\partial_y\!\left(-\frac{\mu_2z}{r_{2j}^3}\right)&\partial_z\!\left(-\frac{\mu_2z}{r_{2j}^3}\right)
\end{bmatrix}
\begin{bmatrix}\hat c_1\\ \hat c_2\\ \hat c_3\end{bmatrix},\qquad r_{2j}=\sqrt{x^2+y^2+z^2}.
\end{equation}
Using $\partial r_{2j}/\partial x=x/r_{2j}$:
\begin{equation}
D_{11}=\frac{\partial}{\partial x}\left(-\frac{\mu_2x}{(x^2+y^2+z^2)^{3/2}}\right)=-\frac{\mu_2}{r_{12}^3}+\frac{3\mu_2x^2}{r_{12}^5},\qquad
D_{12}=\frac{\partial}{\partial y}\left(-\frac{\mu_2x}{(x^2+y^2+z^2)^{3/2}}\right)=\frac{3\mu_2xy}{r_{12}^5},
\end{equation}
and in general
\begin{equation}
D_{ij}=-\frac{\mu_2}{r_{12}^3}\,\delta_{ij}+\frac{3\mu_2\,x_ix_j}{r_{12}^5}.
\end{equation}
The derivatives are evaluated at $\underline r_{21}$; since each $D_{ij}$ is even in the coordinates and $r_{21}=r_{12}$, the result is the same
at $\underline r_{21}=-\underline r_{12}$. In coordinate-free form:
\begin{equation}
\boxed{\;\underrightarrow{D}=\left.\frac{\partial}{\partial\underline r_{2j}}\left[-\frac{\mu_2}{r_{2j}^3}\underline r_{2j}\right]\right|_{\underline r_{21}}
=\frac{\mu_{2}}{r_{12}^{5}}\left[3\,\underline{r}_{12}\underline{r}_{12}-r_{12}^{2}\,\underrightarrow{1}\right]\;}
\qquad \underrightarrow{1}=\hat c_1\hat c_1+\hat c_2\hat c_2+\hat c_3\hat c_3 ,
\end{equation}
where $\underline r_{12}\underline r_{12}$ ($=\underline r_{12}\underline r_{12}^T$) is the outer product. Therefore
\begin{equation}
\underline f_{3B}\simeq\frac{\mu_2}{r_{12}^5}\left[3\,\underline r_{12}\,(\underline r_{12}\cdot\underline r_{1S})-r_{12}^2\,\underline r_{1S}\right]
=\frac{\mu_2}{r_{12}^3}\left[3\,\hat r_{12}(\hat r_{12}\cdot\underline r_{1S})-\underline r_{1S}\right].
\end{equation}

\paragraph{Interpretation.} The perturbation is linear in $r_{1S}$ and proportional to $\mu_2/r_{12}^3$ (tidal acceleration). In a frame with
the first axis along $\hat r_{12}$ the dyad is $\frac{\mu_2}{r_{12}^3}\mathrm{diag}(2,-1,-1)$: the spacecraft is pulled away from the Earth along
the Earth--third-body line (stretching, $+2$) and pushed towards it in the transverse directions (compression, $-1$). Because of
the $1/r_{12}^3$ factor the Moon, though much less massive, has a tidal effect about twice that of the Sun \notinnotes. This form is
used to compute the average effect on $\underline h$ (\qref{C7}).

\begin{keyf}
Condition $r_{1S}\ll r_{21}$; \ $\underline f_{3B}\simeq\dfrac{\partial\underline g}{\partial\underline x}\Big|_{\underline r_{21}}\!\cdot\underline r_{1S}$; \
$D_{ij}=-\dfrac{\mu_2}{r^3}\delta_{ij}+\dfrac{3\mu_2x_ix_j}{r^5}$; \ $\underrightarrow{D}=\dfrac{\mu_2}{r_{12}^5}\left[3\underline r_{12}\underline r_{12}-r_{12}^2\underrightarrow{1}\right]$.
\end{keyf}

\pitfalls
\begin{itemize}
\item State the condition $|\underline r_{1S}|\ll|\underline r_{21}|$ before expanding: it is what the question asks explicitly.
\item The zero-order term $\underline g(\underline r_{21})$ cancels with the second term of $\underline f_{3B}$: only the gradient survives.
\item Show at least one diagonal and one off-diagonal derivative, then the compact form with the unit dyad.
\end{itemize}

\source{cap08.tex l.1211--1329 (approximate perturbing acceleration using the gravity gradient; the dyad and its compact form).}

% ---------------------------------------------------------------------
\question{C9}{Aerodynamic drag: acceleration and effect on the orbit elements}
% source: cap08.tex l.398-559
\begin{qbox}
Report the expression of the drag acceleration and describe its effect on two specific orbit elements.
\qmeta{Jul 2025; Mar 2025 ("aerodynamic drag").}
\end{qbox}

\begin{outline}
\begin{enumerate}
\item Where it matters (100--1000\,km), exponential atmosphere $\rho(\tilde h)$.
\item $\underline a_D=-\frac12C_D\frac{S}{m}\rho v_R^2\hat v_R$, meaning of each term; $\underline v_R=\underline v-\underline\omega_E\times\underline r$; ballistic coefficient $C_DS/m$.
\item Effect: dissipative, along $-\hat v$ ($f_\theta<0$, $f_h\simeq0$) $\Rightarrow$ $a\downarrow$ and $e\downarrow$ (Gauss equations).
\item Elliptic orbit: drag concentrated at perigee $\Rightarrow$ apogee decreases, circularisation; circular orbit: slow spiral.
\item Near-circular: $\dot a=-C_D\frac Sm\rho\sqrt{\mu a}$, $\sqrt{a_f}-\sqrt{a_i}=-\frac{C_DS}{2m}\rho\sqrt\mu\,\Delta t$. Plane nearly unchanged.
\end{enumerate}
\end{outline}

\answer
\paragraph{When drag matters.} By convention space begins at 100\,km (over 99.9999\% of the atmosphere is below it), but up to about
1000\,km the residual atmosphere still perturbs the orbit, and the effect grows rapidly as the altitude decreases. The density is
modelled with a piecewise exponential law (data from USSA76):
\begin{equation}
\rho(\tilde{h})=\rho_{R,i}\exp\!\left[-\frac{\tilde{h}-\tilde{h}_{R,i}}{H_{R,i}}\right],
\end{equation}
$\tilde h_{R,i}$ reference altitude with density $\rho_{R,i}$, $H_{R,i}$ scale height.

\paragraph{Drag acceleration.}
\begin{equation}
\boxed{\;\underline{a}_{D}=\frac{\underline{D}}{m}=-\frac{1}{2}\,C_{D}\,\frac{S}{m}\,\rho\,v_{R}^{2}\,\hat{v}_{R}\;}
\end{equation}
\begin{itemize}
\item $S$ = aerodynamic reference surface (depends on geometry and attitude); $m$ = spacecraft mass (varies if propulsion is used);
\item $C_D$ = drag coefficient: in the rarefied (free-molecular) regime $C_D\simeq2.2$ (minimum 2 for a sphere);
\item $\rho$ = atmospheric density at the current altitude;
\item $\underline v_R=\underline v-\underline v_a$ = velocity relative to the atmosphere, which is assumed to rotate with the Earth:
$\underline v_a=\underline\omega_E\times\underline r=\omega_Er\cos\phi\,\hat E$ (small compared with the orbital speed).
\end{itemize}
$C_D\,S/m$ is the \textbf{ballistic coefficient}: the larger it is, the stronger the action of drag. Drag is \textbf{dissipative}
(non-conservative, a surface perturbation) and opposite to the relative velocity: essentially $f_\theta<0$, small $f_r$, and
$f_h\simeq0$.

\paragraph{Effect on the elements: $a$ and $e$ decrease.} From the Gauss equations
$\dot a=\frac{2a^2}{h}\left(e\,s_{\theta_*}f_r+\frac pr f_\theta\right)$ and
$\dot e=\sqrt{p/\mu}\left[f_rs_{\theta_*}+f_\theta\frac{e+ec_{\theta_*}^2+2c_{\theta_*}}{1+ec_{\theta_*}}\right]$: with $f_\theta<0$ we get
$\dot a<0$; near perigee ($\theta_*\simeq0$) $\dot e\simeq2\sqrt{p/\mu}\,f_\theta<0$.
\begin{itemize}
\item \textbf{Elliptic orbit}: the density is much higher at perigee, so the drag action is \textbf{concentrated at perigee}. Like a small
retro-burn at perigee, it lowers the \textbf{apogee} altitude, while the perigee altitude is only marginally reduced. Hence
\textbf{$a\downarrow$ (semi-major axis) and $e\downarrow$ (eccentricity)}: the orbit is progressively circularised.
\item \textbf{Circular LEO}: the long-term effect is a \textbf{spiral trajectory slowly decreasing in altitude}, until re-entry.
\item The orbit plane is practically unaffected: $f_h^{(D)}$ is very small, so $\dot i\simeq0$, $\dot\Omega\simeq0$.
\end{itemize}
\begin{center}\includegraphics[width=0.45\linewidth]{drag}\end{center}

\paragraph{Approximate analysis for near-circular orbits.} With $a\simeq R$, $v\simeq\sqrt{\mu/a}\,\hat\theta$ and $\underline v_R\simeq\underline v$:
$f_\theta^{(D)}=-\frac12C_D\frac Sm\rho\frac\mu a$, $f_r=f_h=0$, hence
\begin{equation}
\dot{a}=\frac{2a^{3/2}}{\sqrt{\mu}}f_{\theta}^{(D)}=-C_{D}\frac{S}{m}\rho\sqrt{\mu a}
\;\xrightarrow{\ \rho\simeq\text{const}\ }\;
a_{\text{fin}}^{1/2}-a_{\text{ini}}^{1/2}=-\frac{C_{D}S}{2m}\rho\sqrt{\mu}\,(t_{\text{fin}}-t_{\text{ini}}),
\end{equation}
valid for small variations of $a$ (otherwise $\rho$ changes). Note that as $a$ decreases the orbital speed $\sqrt{\mu/a}$ \emph{increases}:
drag removes energy but the satellite speeds up while spiralling down \notinnotes.

\begin{keyf}
$\underline a_D=-\frac12C_D\dfrac Sm\rho v_R^2\hat v_R$, \ $\underline v_R=\underline v-\underline\omega_E\times\underline r$, \ ballistic coeff.\ $C_DS/m$, \
$a\downarrow$, $e\downarrow$; \ $\dot a=-C_D\dfrac Sm\rho\sqrt{\mu a}$.
\end{keyf}

\pitfalls
\begin{itemize}
\item Use the velocity \emph{relative to the atmosphere} and the square of its magnitude.
\item The two elements are $a$ and $e$: explain the perigee mechanism (apogee lowers $\Rightarrow$ circularisation).
\item Drag hardly changes $i$ and $\Omega$.
\end{itemize}

\source{cap08.tex l.398--426 (atmosphere model), l.428--499 (drag acceleration), l.501--519 (effect on elements), l.521--559
(near-circular orbits; errata \#11); Gauss equations l.338--372. Formulario: \fml{8-DRG}, \fml{8-DRC}.}
